\begin{lemma}[Arbitrary interpolation on one block of a pair]
    \label{lem:pair_product_span_arbitrary_first_value}
    \lean{MPSTensor.exists_mem_span_wordTuple_eq_and_eq_zero_of_pairWordTupleSpanTop}
    \leanok
    \uses{def:pair_trace_separation_words}
    Let $(A^\ell)_{\ell=1}^r$ be a finite family of tensors. Suppose that the
    length-$S$ pair evaluations of blocks $k$ and $j$ span their product
    matrix algebra. For every $X\in\MN{D_k}$, there is a tuple
    $M=(M_\ell)_\ell$ in the span of the simultaneous length-$S$ word
    evaluations such that $M_k=X$ and $M_j=0$.
\end{lemma}

\begin{proof}
    \leanok
    \uses{def:pair_trace_separation_words}
    Full pair span gives coefficients satisfying
    $(X,0) = \sum_{|w|=S}c_w(A^k_w,A^j_w)$.
    Define $M_\ell=\sum_{|w|=S}c_wA^\ell_w$ for every block $\ell$. Then $M$
    lies in the simultaneous word-tuple span, and its $k$- and $j$-components
    are $X$ and $0$, respectively.
\end{proof}

\begin{lemma}[Blockwise matrix equality from trace pairings]
    \label{lem:block_matrices_equal_from_trace_pairings}
    \lean{MPSTensor.block_matrices_eq_of_wordTupleSpanTop_trace}
    \leanok
    \uses{def:blockwise_product_word_span}
    Let $A^1,\ldots,A^g$ be block tensors whose length-$m$ simultaneous word
    tuples span the product algebra. If, for every word $w$ of length $m$,
    two blockwise matrix families $X_j$ and $Y_j$ satisfy
    $\sum_j\tr(X_jA^j_w) = \sum_j\tr(Y_jA^j_w)$,
    then $X_j=Y_j$ for every block $j$.
\end{lemma}

\begin{proof}
    \leanok
    \uses{def:blockwise_product_word_span}
    For every word $w$ of length $m$,
    \begin{align}
        \sum_j\tr\!\left((X_j-Y_j)A^j_w\right)
         & =
        \sum_j\tr(X_jA^j_w)
        -
        \sum_j\tr(Y_jA^j_w) =
        0.
        \notag
    \end{align}
    Nondegeneracy of the product trace pairing, together with the spanning
    hypothesis, gives $X_j-Y_j=0$ for every block $j$.
\end{proof}

\begin{lemma}[Boundary identities from block-sum membership]
    \label{lem:blockwise_boundary_identities_from_left_trace_membership}
    \lean{MPSTensor.pgvwc07_boundary_identities_of_leftBoundaryComponent_mem_iSup}
    \leanok
    \uses{lem:block_matrices_equal_from_trace_pairings,
        def:blockwise_product_word_span,
        def:ground_space_parent}
    Let $A^1,\ldots,A^g$ be block tensors whose length-$m$ simultaneous word
    tuples span the product algebra. Let $C^j_a\in \MN{D_j}$ be matrices
    indexed by blocks $j$ and physical letters $a$. Suppose that, for every
    pair of boundary letters $a,b$ and every word $w$ of length $m$, a vector
    $\psi$ has the left-boundary trace form
    $\psi(a,w,b) = \sum_j\tr\!\left(A^j_b C^j_a A^j_w\right)$.
    If $\psi\in\bigvee_j G_{m+2}(A^j)$, then there are matrices $E_j$ such
    that, for every block $j$ and all $a,b$,
    $A^j_bC^j_a = A^j_bE_jA^j_a$.
\end{lemma}

\begin{proof}
    \leanok
    \uses{lem:block_matrices_equal_from_trace_pairings,
        def:ground_space_parent,
        def:ground_space_map}
    Write $\psi=\sum_j\phi_j$ with $\phi_j\in G_{m+2}(A^j)$. For each $j$,
    choose $E_j$ whose image under the length-$m+2$ ground-space map for
    $A^j$ is $\phi_j$. Comparing the coefficient of the word $awb$ in the
    two expressions for $\psi$ gives
    \begin{align}
        \sum_j\tr\!\left(A^j_b C^j_a A^j_w\right)
         & =
        \sum_j\tr\!\left(A^j_bE_jA^j_aA^j_w\right).
        \notag
    \end{align}
    Lemma~\ref{lem:block_matrices_equal_from_trace_pairings} gives the
    displayed blockwise identities.
\end{proof}

\begin{lemma}[Blockwise boundary compatibility from traces]
    \label{lem:blockwise_boundary_compatibility_from_traces}
    \lean{MPSTensor.pgvwc07_blockwise_compatibility_of_trace_decomposition}
    \leanok
    \uses{lem:block_matrices_equal_from_trace_pairings}
    Let $A^1,\ldots,A^g$ be block tensors with the following common
    word-span property: for a fixed $m$, the simultaneous tuples
    $w \longmapsto \left(A^1_w,\ldots,A^g_w\right), \qquad |w|=m$,
    span the full product algebra $\prod_j \MN{D_j}$. Suppose that matrices
    $C^j_a,D^j_b\in \MN{D_j}$ satisfy, for all physical indices $a,b$ and
    all words $w$ of length $m$,
    \begin{align}
        \sum_j \tr\!\left(A^j_b C^j_a A^j_w\right)
         & =
        \sum_j \tr\!\left(D^j_b A^j_a A^j_w\right).
        \notag
    \end{align}
    Then
    $A^j_bC^j_a = D^j_bA^j_a$
    for every block $j$ and all $a,b$.
\end{lemma}

\begin{proof}
    \leanok
    \uses{lem:block_matrices_equal_from_trace_pairings}
    For fixed $a,b$, set $\Delta_j=A^j_bC^j_a-D^j_bA^j_a$. The assumed trace
    equality says that
    $\sum_j \tr\!\left(\Delta_jA^j_w\right) = 0$
    for every word $w$ of length $m$. Since the tuples
    $(A^1_w,\ldots,A^g_w)$ span the product algebra, the product trace
    pairing is nondegenerate and hence every $\Delta_j$ is zero.
\end{proof}

\begin{lemma}[Word-valued boundary compatibility from traces]
    \label{lem:blockwise_word_boundary_compatibility_from_traces}
    \lean{MPSTensor.pgvwc07_blockwise_word_compatibility_of_trace_decomposition,
        MPSTensor.pgvwc07_complementary_word_compatibility_of_trace_decomposition}
    \leanok
    \uses{lem:block_matrices_equal_from_trace_pairings}
    Let $A^1,\ldots,A^g$ be block tensors with the common word-span property
    at length $m$. Suppose that $C^j_\rho\in\MN{D_j}$ is indexed by a block
    $j$ and a word $\rho$ of some fixed length $M$, and that
    $D^j_\beta\in\MN{D_j}$ is indexed by a block $j$ and a word $\beta$ of
    some fixed length $K$. If, for every word $\rho$ of length $M$, every
    word $\beta$ of length $K$, and every middle word $w$ of length $m$,
    \begin{align}
        \sum_j\tr\!\left(A^j_\beta C^j_\rho A^j_w\right)
         & =
        \sum_j\tr\!\left(D^j_\beta A^j_\rho A^j_w\right),
        \notag
    \end{align}
    then
    $A^j_\beta C^j_\rho = D^j_\beta A^j_\rho$
    for every block $j$ and all words $\beta,\rho$. In particular, taking
    $D^j_\beta=X_jA^j_\beta$ gives
    $A^j_\beta C^j_\rho = (X_jA^j_\beta)A^j_\rho$.
\end{lemma}

\begin{proof}
    \leanok
    \uses{lem:block_matrices_equal_from_trace_pairings}
    Fix $\beta$ and $\rho$, and set
    $\Delta_j=A^j_\beta C^j_\rho-D^j_\beta A^j_\rho$. The trace equality
    says that
    $\sum_j\tr(\Delta_jA^j_w) = 0$
    for every word $w$ of length $m$. The same nondegenerate product
    trace-pairing argument as in
    Lemma~\ref{lem:blockwise_boundary_compatibility_from_traces} gives
    $\Delta_j=0$ for every $j$. The specialization
    $D^j_\beta=X_jA^j_\beta$ is immediate.
\end{proof}

\begin{lemma}[Fixed complementary-word compatibility from traces]
    \label{lem:fixed_complementary_word_boundary_compatibility_from_traces}
    \lean{MPSTensor.pgvwc07_fixed_complementary_word_compatibility_of_trace_decomposition}
    \leanok
    \uses{lem:block_matrices_equal_from_trace_pairings}
    Let $A^1,\ldots,A^g$ be block tensors with the common word-span property
    at length $m$. Fix a complementary word $\rho$ of length $M$ and matrices
    $C^j_\rho,X_j\in\MN{D_j}$. If, for every wrapped word $\beta$ of length
    $K$ and every middle word $w$ of length $m$,
    \begin{align}
        \sum_j\tr\!\left(A^j_\beta C^j_\rho A^j_w\right)
         & =
        \sum_j\tr\!\left((X_jA^j_\beta)A^j_\rho A^j_w\right),
        \notag
    \end{align}
    then
    $A^j_\beta C^j_\rho = (X_jA^j_\beta)A^j_\rho$
    for every block $j$ and every wrapped word $\beta$.
\end{lemma}

\begin{proof}
    \leanok
    \uses{lem:block_matrices_equal_from_trace_pairings}
    Fix $\beta$ and set
    $\Delta_j=A^j_\beta C^j_\rho-(X_jA^j_\beta)A^j_\rho$. The trace equality
    gives
    $\sum_j\tr(\Delta_jA^j_w) = 0$
    for every word $w$ of length $m$. The common word-span property and the
    product trace pairing imply $\Delta_j=0$ for every $j$.
\end{proof}

\begin{lemma}[Normalized boundary matrix identities]
    \label{lem:normalized_boundary_matrix_identities}
    \lean{MPSTensor.pgvwc07_boundary_matrix_identities_of_indexed_compatibility}
    \lean{MPSTensor.pgvwc07_boundary_matrix_identities_of_compatibility}
    \lean{MPSTensor.pgvwc07_boundary_word_matrix_identities_of_compatibility}
    \lean{MPSTensor.pgvwc07_boundary_matrix_identities_of_two_index_compatibility}
    \lean{MPSTensor.pgvwc07_boundary_word_matrix_identities_of_two_length_compatibility}
    \leanok
    Let $(A_a)_{a\in I}$, $(C_a)_{a\in I}$, and $(D_b)_{b\in I}$ be matrices
    in $\MD$, indexed by a finite alphabet. If
    \begin{align}
        \sum_a A_a A_a^\dagger & = \Id,
        \notag                             \\
        A_b C_a                & = D_b A_a
        \notag
    \end{align}
    for all $a,b\in I$, and if $E=\sum_a C_a A_a^\dagger$, then
    \begin{align}
        D_b    & = A_bE,
        \notag             \\
        A_bC_a & = A_bEA_a
        \notag
    \end{align}
    for all $a,b\in I$.

    In particular, the same calculation applies when the alphabet is the finite
    set of words of a fixed length and $A_\beta$ denotes
    $A_{\beta_1}\cdots A_{\beta_K}$.

    The same argument also applies when the normalized family is indexed by a
    finite set $J$, while $(A_b)_{b\in I}$ and $(D_b)_{b\in I}$ are indexed
    by a set $I$. If $(B_a)_{a\in J}$ and $(C_a)_{a\in J}$ are indexed by
    $J$, and if, for all $a\in J$ and $b\in I$,
    \begin{align}
        \sum_a B_aB_a^\dagger & = \Id,
        \notag                            \\
        A_bC_a                & = D_bB_a,
        \notag
    \end{align}
    then, with $E=\sum_a C_aB_a^\dagger$, one has
    \begin{align}
        D_b    & = A_bE,
        \notag              \\
        A_bC_a & = A_bEB_a.
        \notag
    \end{align}
    This includes the case where $I$ and $J$ are word sets of different
    lengths.
\end{lemma}

\begin{proof}
    \leanok
    Multiply $D_b$ by the normalization:
    \begin{align}
        D_b
         & =
        D_b\sum_a A_aA_a^\dagger =
        \sum_a D_bA_aA_a^\dagger =
        \sum_a A_bC_aA_a^\dagger =
        A_bE.
        \notag
    \end{align}
    Substituting this identity into $A_bC_a=D_bA_a$ gives
    $A_bC_a=A_bEA_a$. For two alphabets, the corresponding step is
    \begin{align}
        D_b
         & =
        D_b\sum_a B_aB_a^\dagger =
        \sum_a D_bB_aB_a^\dagger =
        \sum_a A_bC_aB_a^\dagger =
        A_bE,
        \notag
    \end{align}
    where $E=\sum_a C_aB_a^\dagger$. Substituting this identity into
    $A_bC_a=D_bB_a$ gives $A_bC_a=A_bEB_a$.
\end{proof}

\begin{lemma}[Explicit $C^j,D^j,E^j$ identities from trace decompositions]
    \label{lem:pgvwc_cde_identities_from_trace_decompositions}
    \lean{MPSTensor.pgvwc07_complementary_word_cde_identities_of_trace_decomposition}
    \leanok
    \uses{lem:sum_evalWord_mul_conjTranspose_evalWord,
        lem:blockwise_word_boundary_compatibility_from_traces,
        lem:normalized_boundary_matrix_identities}
    Let $A^1,\ldots,A^g$ be block tensors with the common word-span property
    at length $m$. Fix matrices $X_j$, and put
    $D^j_\beta=X_jA^j_\beta$. Suppose that $C^j_\rho$ is indexed by a block
    $j$ and a complementary word $\rho$ of length $M$, and assume that, for
    every wrapped word $\beta$ of length $K$, every complementary word
    $\rho$, and every middle word $w$ of length $m$,
    \begin{align}
        \sum_j\tr\!\left(A^j_\beta C^j_\rho A^j_w\right)
         & =
        \sum_j\tr\!\left(D^j_\beta A^j_\rho A^j_w\right).
        \notag
    \end{align}
    If each block satisfies
    $\sum_a A^j_aA^{j\,\dagger}_a = \Id$,
    then, for every block $j$, there is a matrix
    $E^j = \sum_\rho C^j_\rho\left(A^j_\rho\right)^\dagger$
    such that, for every wrapped word $\beta$ and complementary word $\rho$,
    \begin{align}
        D^j_\beta          & = A^j_\beta E^j,
        \notag                                        \\
        A^j_\beta C^j_\rho & = A^j_\beta E^jA^j_\rho.
        \notag
    \end{align}
    This is the calculation in
    \cite[Theorem~12, proof lines~1446--1451]{PerezGarcia2007Matrix}, with
    the adjoint on $A^j_\rho$ required by the displayed normalization.
\end{lemma}

\begin{proof}
    \leanok
    \uses{lem:sum_evalWord_mul_conjTranspose_evalWord,
        lem:blockwise_word_boundary_compatibility_from_traces,
        lem:normalized_boundary_matrix_identities}
    Lemma~\ref{lem:blockwise_word_boundary_compatibility_from_traces} gives
    $A^j_\beta C^j_\rho = D^j_\beta A^j_\rho$.
    The single-letter normalization gives the word normalization
    $\sum_\rho A^j_\rho\left(A^j_\rho\right)^\dagger = \Id$.
    Lemma~\ref{lem:normalized_boundary_matrix_identities}, applied to the
    two word alphabets, gives
    \begin{align}
        E^j
         & =
        \sum_\rho C^j_\rho\left(A^j_\rho\right)^\dagger,
        \notag \\
        D^j_\beta
         & =
        A^j_\beta E^j,
        \notag \\
        A^j_\beta C^j_\rho
         & =
        A^j_\beta E^jA^j_\rho.
        \notag
    \end{align}
\end{proof}

\begin{lemma}[Block-diagonal boundary $C^j,D^j,E^j$ identities from trace decompositions]
    \label{lem:block_boundary_cde_identities_from_trace_decompositions}
    \lean{MPSTensor.pgvwc07_complementary_word_cde_identities_of_block_boundary_trace_decomposition}
    \leanok
    \uses{lem:pgvwc_cde_identities_from_trace_decompositions}
    Let $A^1,\ldots,A^g$ be block tensors with the common word-span property
    at length $m$. Fix $N\in\N$, complex numbers $\mu_j$, and block-diagonal
    boundary matrices $X_j$, and put
    $D^j_\beta=(\mu_j^NX_j)A^j_\beta$. Suppose that, for every wrapped word
    $\beta$ of length $K$, every complementary word $\rho$ of length $M$,
    and every middle word $w$ of length $m$,
    \begin{align}
        \sum_j\tr\!\left(A^j_\beta C^j_\rho A^j_w\right)
         & =
        \sum_j\tr\!\left(((\mu_j^NX_j)A^j_\beta)A^j_\rho A^j_w\right).
        \notag
    \end{align}
    If each block satisfies
    $\sum_a A^j_aA^{j\,\dagger}_a = \Id$,
    then, for every block $j$, there is a matrix
    $E^j = \sum_\rho C^j_\rho\left(A^j_\rho\right)^\dagger$
    such that, for every wrapped word $\beta$ and complementary word $\rho$,
    \begin{align}
        (\mu_j^NX_j)A^j_\beta & = A^j_\beta E^j,
        \notag                                           \\
        A^j_\beta C^j_\rho    & = A^j_\beta E^jA^j_\rho.
        \notag
    \end{align}
    This is the boundary-matrix calculation in
    \cite[Theorem~12, proof lines~1446--1451]{PerezGarcia2007Matrix},
    after substituting the block-diagonal boundary matrix
    $D^j_\beta=(\mu_j^NX_j)A^j_\beta$.
\end{lemma}

\begin{proof}
    \leanok
    \uses{lem:pgvwc_cde_identities_from_trace_decompositions}
    Apply Lemma~\ref{lem:pgvwc_cde_identities_from_trace_decompositions}
    with $D^j_\beta=(\mu_j^NX_j)A^j_\beta$.
\end{proof}

\begin{lemma}[Complementary-word boundary identities from a compatibility hypothesis]
    \label{lem:complementary_word_boundary_identities_from_pgvwc}
    \lean{MPSTensor.pgvwc07_complementary_word_boundary_identities_formula_of_compatibility}
    \leanok
    \uses{lem:normalized_boundary_matrix_identities}
    Let $A_\beta$ be the products over wrapped words of length $K$, let
    $A_\rho$ be the products over complementary words of length $M$, and let
    $X\in\MD$ be a matrix. Suppose that the complementary-word products
    satisfy
    $\sum_\rho A_\rho A_\rho^\dagger = \Id$
    and that there are matrices $C_\rho$, indexed by complementary words,
    such that, for every wrapped word $\beta$ and complementary word $\rho$,
    $A_\beta C_\rho = (X A_\beta)A_\rho$.
    Then, for every complementary word $\rho$, with
    \begin{align}
        E_\rho
         & =
        \left(\sum_\gamma C_\gamma A_\gamma^\dagger\right)A_\rho,
        \notag
    \end{align}
    one has, for every wrapped word $\beta$,
    $(X A_\beta)A_\rho = A_\beta E_\rho$.
\end{lemma}

\begin{proof}
    \leanok
    \uses{lem:normalized_boundary_matrix_identities}
    Apply Lemma~\ref{lem:normalized_boundary_matrix_identities} with
    $D_\beta=X A_\beta$. It gives
    $A_\beta C_\rho = A_\beta E A_\rho$
    for $E=\sum_\gamma C_\gamma A_\gamma^\dagger$. With $E_\rho=EA_\rho$,
    the compatibility hypothesis gives
    \begin{align}
        (X A_\beta)A_\rho
         & =
        A_\beta C_\rho =
        A_\beta E A_\rho =
        A_\beta E_\rho.
        \notag
    \end{align}
\end{proof}

\begin{lemma}[Complementary-word boundary identities from trace decompositions]
    \label{lem:complementary_word_boundary_identities_from_trace_decompositions}
    \lean{MPSTensor.pgvwc07_complementary_word_boundary_identities_of_trace_decomposition}
    \leanok
    \uses{lem:sum_evalWord_mul_conjTranspose_evalWord,
        lem:blockwise_word_boundary_compatibility_from_traces,
        lem:complementary_word_boundary_identities_from_pgvwc}
    Let $A^1,\ldots,A^g$ be block tensors with the common word-span property
    at length $m$. Fix matrices $X_j$. Suppose that $C^j_\rho$ is indexed
    by a block $j$ and a complementary word $\rho$ of length $M$. Assume
    that, for every wrapped word $\beta$ of length $K$, every complementary
    word $\rho$, and every middle word $w$ of length $m$,
    \begin{align}
        \sum_j\tr\!\left(A^j_\beta C^j_\rho A^j_w\right)
         & =
        \sum_j\tr\!\left((X_jA^j_\beta)A^j_\rho A^j_w\right).
        \notag
    \end{align}
    If each block satisfies
    $\sum_a A^j_aA^{j\,\dagger}_a = \Id$,
    then, for every block $j$ and every complementary word $\rho$, there is
    a matrix $E_{j,\rho}$ such that, for every wrapped word $\beta$,
    $(X_jA^j_\beta)A^j_\rho = A^j_\beta E_{j,\rho}$.
\end{lemma}

\begin{proof}
    \leanok
    \uses{lem:sum_evalWord_mul_conjTranspose_evalWord,
        lem:blockwise_word_boundary_compatibility_from_traces,
        lem:complementary_word_boundary_identities_from_pgvwc}
    Lemma~\ref{lem:blockwise_word_boundary_compatibility_from_traces}, applied
    with $D^j_\beta=X_jA^j_\beta$, gives
    $A^j_\beta C^j_\rho = (X_jA^j_\beta)A^j_\rho$.
    The hypothesis $\sum_a A^j_a A^{j\,\dagger}_a=\Id$ extends to words of
    length $M$:
    $\sum_\rho A^j_\rho\left(A^j_\rho\right)^\dagger = \Id$.
    Lemma~\ref{lem:complementary_word_boundary_identities_from_pgvwc} gives
    the matrices $E_{j,\rho}$.
\end{proof}

\begin{lemma}[A left-boundary summand is a ground-space vector]
    \label{lem:left_boundary_component_mem_ground_space}
    \lean{MPSTensor.pgvwc07LeftBoundaryComponent_mem_groundSpace}
    \leanok
    Let $A$ be a block tensor, and let $C_a,E\in\MD$ satisfy
    $A_b C_a = A_bEA_a$
    for all physical indices $a,b$. For a word
    $\sigma=(\sigma_1,\ldots,\sigma_{n+2})$, define
    \begin{align}
        \alpha(\sigma)
         & =
        \tr\!\left(A_{\sigma_{n+2}} C_{\sigma_1}
        A_{\sigma_2}\cdots A_{\sigma_{n+1}}\right).
        \notag
    \end{align}
    Then $\alpha\in G_{n+2}(A)$.
\end{lemma}

\begin{proof}
    \leanok
    The boundary identity gives
    \begin{align}
        A_{\sigma_{n+2}} C_{\sigma_1}
         & =
        A_{\sigma_{n+2}}EA_{\sigma_1}.
        \notag
    \end{align}
    Hence, by cyclicity of the trace,
    \begin{align}
        \alpha(\sigma)
         & =
        \tr\!\left(A_{\sigma_1}A_{\sigma_2}\cdots A_{\sigma_{n+2}}E\right),
        \notag
    \end{align}
    which is the ground-space parametrization of $E$.
\end{proof}

\begin{lemma}[Left-boundary summands are ground-space vectors]
    \label{lem:left_boundary_summands_mem_block_ground_spaces}
    \lean{MPSTensor.pgvwc07_sum_leftBoundaryComponents_mem_iSup_groundSpace}
    \leanok
    \uses{lem:left_boundary_component_mem_ground_space}
    Let $A^1,\ldots,A^g$ be block tensors, and let
    $C^j_a,E^j\in\MN{D_j}$ satisfy
    $A^j_b C^j_a = A^j_bE^jA^j_a$
    for every block $j$ and all physical indices $a,b$. For a word
    $\sigma=(\sigma_1,\ldots,\sigma_{n+2})$, define
    \begin{align}
        \alpha_j(\sigma)
         & =
        \tr\!\left(A^j_{\sigma_{n+2}} C^j_{\sigma_1}
        A^j_{\sigma_2}\cdots A^j_{\sigma_{n+1}}\right).
        \notag
    \end{align}
    Then
    $\sum_j \alpha_j \in \bigvee_j G_{n+2}(A^j)$.
\end{lemma}

\begin{proof}
    \leanok
    \uses{lem:left_boundary_component_mem_ground_space}
    Fix $j$. By the assumed boundary identity and cyclicity of the trace,
    \begin{align}
        \tr\!\left(A^j_{\sigma_{n+2}} C^j_{\sigma_1}
        A^j_{\sigma_2}\cdots A^j_{\sigma_{n+1}}\right)
         & =
        \tr\!\left(A^j_{\sigma_1}A^j_{\sigma_2}\cdots
        A^j_{\sigma_{n+2}}E^j\right).
        \notag
    \end{align}
    Thus $\alpha_j$ is the image of $E^j$ under the parametrization of
    $G_{n+2}(A^j)$. Hence each $\alpha_j$ lies in $G_{n+2}(A^j)$, and the
    finite sum lies in the supremum of these subspaces.
\end{proof}

\begin{lemma}[Left-boundary trace decompositions lie in the block ground spaces]
    \label{lem:left_boundary_trace_decomposition_mem_block_ground_spaces}
    \lean{MPSTensor.pgvwc07_mem_iSup_groundSpace_of_trace_decomposition}
    \leanok
    Let $A^1,\ldots,A^g$ be block tensors with common word-span separation at
    length $n$, and suppose that
    $\sum_a A^j_a A^{j\,\dagger}_a = \Id$
    for every block $j$. Suppose a vector
    $\psi\in(\C^d)^{\otimes(n+2)}$ has the left-boundary trace decomposition
    \begin{align}
        \psi
         & =
        \sum_j\alpha_j,
        \notag \\
        \alpha_j(i_1,\ldots,i_{n+2})
         & =
        \tr\!\left(A^j_{i_{n+2}} C^j_{i_1}
        A^j_{i_2}\cdots A^j_{i_{n+1}}\right),
        \notag
    \end{align}
    and suppose that the two trace decompositions agree for every middle word:
    \begin{align}
        \sum_j
        \tr\!\left(A^j_b C^j_a A^j_w\right)
         & =
        \sum_j
        \tr\!\left(D^j_b A^j_a A^j_w\right).
        \notag
    \end{align}
    Then
    $\psi \in \bigvee_j G_{n+2}(A^j)$.
\end{lemma}

\begin{proof}
    \leanok
    The equality of the two boundary trace expressions and the common word-span
    hypothesis give
    $A^j_bC^j_a = D^j_bA^j_a$
    for every $j,a,b$. By the normalization,
    \begin{align}
        D^j_b & = D^j_b\sum_a A^j_aA^{j\,\dagger}_a = \sum_a A^j_bC^j_aA^{j\,\dagger}_a = A^j_bE^j,
        \notag                                                                                      \\
        E^j   & := \sum_a C^j_aA^{j\,\dagger}_a.
        \notag
    \end{align}
    Hence $A^j_bC^j_a=A^j_bE^jA^j_a$.
    Lemma~\ref{lem:left_boundary_summands_mem_block_ground_spaces} then gives
    $\psi = \sum_j\alpha_j \in \bigvee_j G_{n+2}(A^j)$.
\end{proof}
